\documentclass{scrartcl}
\usepackage[utf8]{inputenc}
\usepackage{hyperref}
\usepackage{url}
\usepackage{natbib}
\usepackage{graphicx}
\usepackage{float}
\usepackage{listings}
\usepackage{xcolor}
\usepackage{cleveref} % this must be the last package to be loaded

\title{\LARGE
  Assignment 1: Poool
}
\subtitle{Concurrent and Distributed Programming\\a.y. 2025-2026}

% Consider watching:
% https://www.youtube.com/watch?v=ihxSUsJB_14
% https://www.youtube.com/watch?v=XTFWaV55uDo

\author{
    Bagattoni Federico \and Sbaraccani Pietro \and Venturini Luca
}

\date{\today}

\begin{document}

\maketitle

\section{Analysis}
\subsection{General requirements}
It is requested to develop a game of pool where players control a ball and fight to
score points by pushing other, smaller, balls inside of the corner holes. 

One ball must be controlled by the player and the other controlled by the computer. At the
end of the game, the player with the most points wins.

It is required to develop two solutions: one using default Java multithread programming
and one using the Java Executor Framework. It is also necessary to use high-level constructs
such as \textbf{monitors} to manage concurrency.

\subsection{Analysis of the sequential solution}
The sequential solution, available as \texttt{sketch01}, when run with large
quantities of balls, is not efficient because at every update,
the single thread running the loop has to cycle multiple times throught every
game object to calculate collisions and update movements.
The the loop has to wait for the view to render the frame with the new update,
resulting in a performance of about less than $\sim$20 frames per second (FPS).

A solution adopting concurrent programming could take advantage of all
the available cores of the machine thus improving calculations speed.
Also, the collisions calculations and the view's rendering could be 
approached in a more efficient way.
\subsection{Analysis of the problem}
% as explained in the course, first we identify the tasks, then we map them to 
% active components
\label{analysis:prob}
The critical tasks identified are:
\begin{itemize}
    \item \textbf{movement calculations}: every ball must be moved in a new position according
        to its current speed
    \item \textbf{collision calculations}: every ball colliding with other balls must change
        direction and speed accordingly
    \item \textbf{frame rendering}: the updated game state is passed to the view that renders it.
        It's important to wait for the complete rendering of the frame before the next update.
\end{itemize}

\section{Design}
\subsection{Architecture}
The architecture adopted is the \textbf{Model-View-Controller}. The controller is composed
of all the threads needed for the computation, in a \textbf{master-worker} architecture
where the control loop, a thread itself, acts as master towards the other threads
synchronisiting their actions and giving them tasks to execute.
\subsection{Components}
The \textbf{active components} needed are threads that will execute work on behalf 
of the controller. These components include also the thread rendering the view, in
the Java implementation known as the Event Dispatcher Thread.

The \textbf{passive components} needed are:
\begin{itemize}
    \item data structures that have to be shared between the threads,
        implemented as monitors
    \item barriers, to synchronize the threads' (master and workers) activities during 
        the calculation phases 
\end{itemize}

\subsection{Modeling}
The tasks mentionted in the section \ref{analysis:prob} are modeled as runnable functions
so that the master can send them to the worker who will execute them.

The barriers, one for each computational task, are modeled with objects exposing a
\texttt{hit()} method so that the active components can hit it once their work is done.
A barrier between the View and the Controller is also necessary but implemented so that
the call to update the view blocks the controller until the frame is painted.

Each ball is modeled as a monitor to prevent races during the collisions calculation.

The class that contains all the balls, the Board, is modeled itself as monitor to ensure
that there's not critical races when the workers are operating on it.

Other domain entities that do not participate on concurrent activities
such as speed, position and boundary are implemented according to their
specifications.

\section{Implementation}
For both versions it has been decided to create as many threads as the available 
cores to exploit all the system's resources. In fact creating a larger number 
of threads would mean more context switching and it would prove damaging to
performance.



\section{Verification}

\section{Conclusion}

\end{document}